% MSc dissertation example file, February 2022
%
% Leave one of the documentclass lines uncommented to match your degree.
% You may remove the logo option if it causes problems.
% Do not change any other options.
% \documentclass[logo,msc,adi]{infthesis}     % Adv Design Inf
% \documentclass[logo,msc,ai]{infthesis}      % AI
% \documentclass[logo,msc,cogsci]{infthesis}  % Cognitive Sci
% \documentclass[logo,msc,cs]{infthesis}      % Computer Sci
% \documentclass[logo,msc,cyber]{infthesis}   % Cyber Sec
% \documentclass[logo,msc,datasci]{infthesis} % Data Sci
% \documentclass[logo,msc,di]{infthesis}      % Design Inf
% \documentclass[logo,msc,dsti]{infthesis}    % Data Sci TI
% \documentclass[logo,msc,inf]{infthesis}     % Informatics
\documentclass[logo,msc]{infthesis}           % degree unspecified, do not change except to add your degree
%%%%%%%%%%%%%%%%%%%%%%%%
% Understand any problems and seek approval before assuming it's ok to remove ugcheck.
\usepackage{msccheck}

% Include any packages you need below, but don't include any that change the page
% layout or style of the dissertation. By including the ugcheck package above,
% you should catch most accidental changes of page layout though.

\usepackage{microtype} % recommended, but you can remove if it causes problems

\begin{document}
\begin{preliminary}

\title{This is the Project Title}

\author{Your Name}

\date{\today}

\abstract{
This skeleton demonstrates how to use the \texttt{infthesis} style for
MSc dissertations in the School of Informatics. It also emphasises the
page limit and associated style restrictions for Informatics dissertations
with course code \texttt{INFR11077}. If your degree has a different project
course code, then it is likely to have different formatting rules.
The file \texttt{skeleton.tex} generates this document and should be used as a
starting point for your thesis. Replace this abstract text with a concise
summary of your report.
}

\maketitle

\newenvironment{ethics}
   {\begin{frontenv}{Research Ethics Approval}{\LARGE}}
   {\end{frontenv}\newpage}

\begin{ethics}
\textbf{Instructions:} \emph{Agree with your supervisor which
statement you need to include. Then delete the statement that you are not using,
and the instructions in italics.\\
\textbf{Either complete and include this statement:}}\\ % DELETE THESE INSTRUCTIONS
%
% IF ETHICS APPROVAL WAS REQUIRED:
This project obtained approval from the Informatics Research Ethics committee.\\
Ethics application number: ???\\
Date when approval was obtained: YYYY-MM-DD\\
%
\emph{[If the project required human participants, edit as appropriate, otherwise delete:]}\\ % DELETE THIS LINE
The participants' information sheet and a consent form are included in the appendix.\\
%
% IF ETHICS APPROVAL WAS NOT REQUIRED:
\textbf{\emph{Or include this statement:}}\\ % DELETE THIS LINE
This project was planned in accordance with the Informatics Research
Ethics policy. It did not involve any aspects that required approval
from the Informatics Research Ethics committee.

\standarddeclaration
\end{ethics}


\begin{acknowledgements}
Any acknowledgements go here.
\end{acknowledgements}


\tableofcontents
\end{preliminary}


\chapter{Introduction}

% Keep the dissertation-format warnings in the file, but leave them commented.
% Do not change the page layout, margins, or line spacing.
% The main body must stay within the page limit.

\section{Motivation}
% Why this forest growth problem matters.
% Why LiDAR-based height prediction is useful.
% Why environmental and spatial effects are worth testing.

\section{Problem Statement and Research Gap}
% What existing approaches do not explain well.
% Why classical growth curves alone are not enough.
% Why you compare DNNs, PINNs, and attribution models.

\section{Aim and Objectives}
% State the main aim.
% List the concrete objectives.
% Keep this short and direct.

\section{Research Questions}
% Prediction question.
% Environmental feature question.
% Spatial attribution question.

\section{Dissertation Structure}
% Brief map of the later chapters.

\chapter{Related Work}

\section{Forest Height and Site Productivity}
% Background on dominant height and site quality.

\section{Classical Forest Growth Modelling}
% Traditional growth models used as baselines.

\paragraph{Yield Class and Site Index}
% Site productivity framing.

\paragraph{Chapman--Richards Curves}
% Main growth-curve baseline.

\paragraph{Site-Specific Growth Curves}
% How curve parameters vary across stands.

\section{LiDAR-Based Forest Monitoring}
% How LiDAR is used to estimate canopy height and change.

\section{Machine Learning for Forest Height Prediction}
% Prior data-driven work.

\paragraph{Statistical and Tree-Based Methods}
% Linear and ensemble baselines.

\paragraph{Deep Neural Networks}
% Why DNNs are a natural comparison.

\section{Physics-Informed Neural Networks}
% Why and how physics constraints are added.

\section{Environmental Drivers of Growth}
% Terrain, climate, wind, soil, and management context.

\section{Spatial Generalisation and Validation}
% Why random splits can overstate performance.

\paragraph{Spatial Autocorrelation}
% Why nearby plots are not independent.

\paragraph{Geographically Varying Models}
% Motivation for local attribution methods.

\section{Research Gap}
% One short summary of the gap your thesis addresses.

\chapter{Data}

\section{Study Area and Data Sources}
% Forest setting, cohorts, and source datasets.

\section{LiDAR and Inventory Data}
% Height target, survey years, and core measured variables.

\section{Environmental and Management Data}
% Terrain, wind, climate, soil, spatial position, and management features.

\section{Data Extraction and Alignment}
% How data were matched to plots and surveys.

\section{Cleaning and Cohort Construction}
% Filtering, exclusions, and 4survey / 6survey construction.

\section{Target Definition and Feature Selection}
% Why top height is the target.
% Why some variables were excluded.
% Why leakage and multicollinearity mattered.

\section{Data Quality and Disturbance Screening}
% Missingness, outliers, and disturbed plots.

\section{Exploratory Data Analysis}
% Patterns that informed later modelling choices.

\chapter{Methodology}

\section{Overview of the Modelling Pipeline}
% Move from baselines to neural models to attribution.

\section{Validation Design}
% Plot-held-out, spatial CV, and temporal splits.

\section{Baseline Models}
% Classical benchmarks for comparison.

\paragraph{Chapman--Richards Baseline}
% Main growth-curve benchmark.

\paragraph{Age-Only and Linear Baselines}
% Simple reference models.

\paragraph{Random Forest Baseline}
% Nonlinear non-neural benchmark.

\section{Deep Neural Networks}
% DNN architecture, inputs, and training.

\paragraph{Age-Only DNN}
% Baseline neural model without environmental inputs.

\paragraph{Environmental DNN}
% DNN with terrain and related features.

\section{Physics-Informed Neural Networks}
% PINN setup and physics constraint.

\paragraph{No-Environment PINN}
% Direct comparison against the age-only DNN.

\paragraph{Environmental PINN}
% PINN with environmental inputs.

\paragraph{Constraint Weight Experiments}
% Test how strongly the physics term should matter.

\section{Environmental Conditioning Variants}
% Newer models that let environmental variables shape the curve.

\section{Growth Curve Attribution}
% Explain variation in fitted growth curves.

\paragraph{Shared-Curve Residual Modelling}
% Pooled residual approach.

\paragraph{Plot-Specific Local Deviation Modelling}
% Local deviations from the shared curve.

\paragraph{Residual Learning Models}
% Elastic Net, XGBoost, and GNNWR.

\section{Evaluation Metrics}
% R squared, residual structure, and spatial diagnostics.

\chapter{Results and Discussion}

\section{Data Summary and Preliminary Findings}
% Main properties of the datasets used.
% Early patterns that motivated the modelling choices.

\section{Baseline Model Results}
% Compare classical baselines across validation settings.

\section{Deep Learning and PINN Results}
% Compare DNN and PINN variants.
% Show where environmental inputs help or hurt.
% Show what the physics term changes.

\section{Robustness and Spatial Error Structure}
% Fold variation, uncertainty, and residual autocorrelation.

\section{Growth Curve Attribution Results}
% Shared-curve and plot-specific attribution results.
% Compare the candidate explanation models.

\section{Interpretation}
% What the results suggest about growth, site effects, and generalisation.

\section{Comparison with Related Work}
% How the findings fit the literature.

\section{Limitations}
% Data limits, modelling limits, and open issues.

\section{Implications for Further Modelling}
% What to test next.

\chapter{Conclusion}

\section{Summary of the Work}
% Brief project summary.

\section{Answers to the Research Questions}
% Short direct answers.

\section{Main Contributions}
% What this thesis adds.

\section{Future Work}
% What remains to be done.

\bibliographystyle{plain}
\bibliography{mybibfile}


% You may delete everything from \appendix up to \end{document} if you don't need it.
\appendix

\chapter{First appendix}

\section{First section}

Any appendices, including any required ethics information, should be included
after the references.

Markers do not have to consider appendices. Make sure that your contributions
are made clear in the main body of the dissertation (within the page limit).

\chapter{Participants' information sheet}

If you had human participants, include key information that they were given in
an appendix, and point to it from the ethics declaration.

\chapter{Participants' consent form}

If you had human participants, include information about how consent was
gathered in an appendix, and point to it from the ethics declaration.


\end{document}
